% ═══════════════════════════════════════════════════════════════════════════════
% EBF APPENDIX QQQ: QUALITY ASSESSMENT FRAMEWORK
% ═══════════════════════════════════════════════════════════════════════════════
% Version: 1.0
% Category: LIT
% Language: English
% Date: 2026-01-06
% Status: Initial Release
% Template Compliance: 100%
% ═══════════════════════════════════════════════════════════════════════════════

\documentclass[11pt,a4paper]{article}
\usepackage[utf8]{inputenc}
\usepackage{amsmath,amssymb,amsthm}
\usepackage{booktabs,longtable,array}
\usepackage{hyperref}
\usepackage{xcolor}
\usepackage{tcolorbox}
\usepackage{enumitem}

% Epistemic Tag Colors
\definecolor{empgreen}{RGB}{34,197,94}
\definecolor{thrblue}{RGB}{59,130,246}
\definecolor{llmyellow}{RGB}{234,179,8}
\definecolor{illorange}{RGB}{249,115,22}
\definecolor{hypred}{RGB}{239,68,68}

% Custom commands
\newcommand{\tagEMP}[1]{\textcolor{empgreen}{\textbf{[EMP:} #1\textbf{]}}}
\newcommand{\tagTHR}[1]{\textcolor{thrblue}{\textbf{[THR:} #1\textbf{]}}}
\newcommand{\tagLLM}[1]{\textcolor{llmyellow}{\textbf{[LLM:} #1\textbf{]}}}
\newcommand{\tagILL}[1]{\textcolor{illorange}{\textbf{[ILL:} #1\textbf{]}}}
\newcommand{\tagHYP}[1]{\textcolor{hypred}{\textbf{[HYP:} #1\textbf{]}}}

\newtheorem{axiom}{Axiom}[section]
\newtheorem{definition}{Definition}[section]
\newtheorem{result}{Result}[section]

\title{\textbf{Appendix UNMAPPED_QLY: Quality Assessment Framework}\\
\large TERAN-Based Content Quality Evaluation for EBF}
\author{EBF Research Group}
\date{Version 1.0 — January 2026}

\begin{document}
\maketitle

% ═══════════════════════════════════════════════════════════════════════════════
% HEADER BLOCK
% ═══════════════════════════════════════════════════════════════════════════════

\begin{tcolorbox}[colback=blue!5,colframe=blue!75!black,title=\textbf{Header Block}]
\begin{tabular}{p{4cm}p{10cm}}
\textbf{META Question:} & HOW GOOD — How do we assess content quality? \\
\textbf{Category:} & META (Framework Infrastructure) \\
\textbf{Scope:} & M1 (Methodological Foundation) \\
\textbf{Dependencies:} & None (foundational) \\
\textbf{Dependents:} & All EBF Appendices \\
\textbf{Output:} & TERAN Score $Q \in [1, 10]$, Epistemic Distribution \\
\end{tabular}

\vspace{0.3cm}
\textbf{Relationship to 5C CORE:}
\begin{center}
\begin{tabular}{|c|c|c|c|c|}
\hline
\textbf{AAA (WHO)} & \textbf{B (HOW)} & \textbf{C (WHAT)} & \textbf{V (WHEN)} & \textbf{BBB (WHERE)} \\
\hline
Assessed & Assessed & Assessed & Assessed & Assessed \\
\hline
\end{tabular}
\end{center}
\end{tcolorbox}

% ═══════════════════════════════════════════════════════════════════════════════
% QUICK REFERENCE BOX
% ═══════════════════════════════════════════════════════════════════════════════

\begin{tcolorbox}[colback=green!5,colframe=green!75!black,title=\textbf{Quick Reference}]

\textbf{The TERAN Formula:}
\begin{equation}
Q = 0.25 \cdot T + 0.30 \cdot E + 0.15 \cdot R + 0.15 \cdot A + 0.15 \cdot N
\end{equation}

\textbf{Five Dimensions:}
\begin{center}
\begin{tabular}{clcl}
\toprule
\textbf{Dim} & \textbf{Name} & \textbf{Weight} & \textbf{Assesses} \\
\midrule
T & Theory & 25\% & Theoretical foundation \\
E & Evidence & 30\% & Empirical grounding \\
R & Rigor & 15\% & Formal precision \\
A & Applicability & 15\% & Practical usefulness \\
N & Novelty/Transparency & 15\% & Epistemic honesty \\
\bottomrule
\end{tabular}
\end{center}

\textbf{Rating Scale:}
\begin{center}
\begin{tabular}{cccl}
\toprule
\textbf{Score} & \textbf{Stars} & \textbf{Label} & \textbf{Meaning} \\
\midrule
9.0--10.0 & $\star\star\star\star\star$ & Publication Ready & Top journal standard \\
7.0--8.9 & $\star\star\star\star\phantom{\star}$ & Minor Revisions & Solid, needs polish \\
5.0--6.9 & $\star\star\star\phantom{\star\star}$ & Working Paper & Substantial work needed \\
3.0--4.9 & $\star\star\phantom{\star\star\star}$ & Draft & Major gaps \\
1.0--2.9 & $\star\phantom{\star\star\star\star}$ & Concept & Placeholder only \\
\bottomrule
\end{tabular}
\end{center}

\textbf{Jump to:} §\ref{sec:foundation} Foundation | §\ref{sec:dimensions} Dimensions | §\ref{sec:tags} Epistemic Tags | §\ref{sec:application} Application
\end{tcolorbox}

% ═══════════════════════════════════════════════════════════════════════════════
% ABSTRACT
% ═══════════════════════════════════════════════════════════════════════════════

\begin{abstract}
This appendix establishes the Quality Assessment Framework for EBF, distinguishing \emph{template compliance} (formal structure) from \emph{content quality} (scientific substance). The TERAN framework—assessing Theory, Evidence, Rigor, Applicability, and Transparency—is derived from the peer-reviewed Quality Assessment Framework (QAF) of Belcher et al.\ (2016), supplemented by the TACT framework (2023), Glassick Criteria (1997), and Guba's naturalistic evaluation criteria (1981). We define five epistemic status tags (EMP, THR, LLM, ILL, HYP) to ensure transparency about the evidential basis of all claims and parameters. The framework enables systematic quality tracking and provides actionable guidance for improving EBF documentation toward publication readiness.
\end{abstract}




\begin{tcolorbox}[
    colback=orange!5!white,
    colframe=orange!75!black,
    title={\textbf{Appendix Scope: Ziel / In-Scope / Out-of-Scope / Constraints}},
    fonttitle=\bfseries
]

\textbf{Ziel (Objective):}
\begin{itemize}[nosep]
    \item Provide systematic analysis for LIT integration
\end{itemize}

\textbf{In-Scope:}
\begin{itemize}[nosep]
    \item Core analysis and findings
    \item Framework integration points
\end{itemize}

\textbf{Out-of-Scope:}
\begin{itemize}[nosep]
    \item Detailed empirical replication $\rightarrow$ METHOD appendices
\end{itemize}

\textbf{Constraints:}
\begin{itemize}[nosep]
    \item Analysis bounded by available data
\end{itemize}

\end{tcolorbox}

%%%%%%%%%%%%%%%%%%%%%%%%%%%%%%%%%%%%%%%%%%%%%%%%%%%%%%%%%%%%%%%%%%%%%%%%%%%%%%%
\section{The Fundamental Question}
\label{sec:qly-fundamental}
%%%%%%%%%%%%%%%%%%%%%%%%%%%%%%%%%%%%%%%%%%%%%%%%%%%%%%%%%%%%%%%%%%%%%%%%%%%%%%%

\begin{tcolorbox}[colback=red!5!white, colframe=red!75!black,
    title=The Fundamental Question]
What are the key mechanisms and implications of this analysis for the
Evidence-Based Framework (EBF)?
\end{tcolorbox}

% ═══════════════════════════════════════════════════════════════════════════════
% SECTION 1: THEORETICAL FOUNDATION
% ═══════════════════════════════════════════════════════════════════════════════

\section{Theoretical Foundation}
\label{sec:foundation}

\subsection{The Quality Gap}

EBF appendices may achieve 100\% \emph{template compliance}—proper headers, cross-references, glossary links—while remaining inadequate in \emph{content quality}. This appendix addresses that gap.

\begin{definition}[Template Compliance vs.\ Content Quality]
\begin{itemize}
    \item \textbf{Template Compliance:} Does the document have the required structural elements?
    \item \textbf{Content Quality:} Is the substance scientifically sound, empirically grounded, and epistemically transparent?
\end{itemize}
\end{definition}

\subsection{Source Frameworks}

The TERAN framework synthesizes five established, peer-reviewed quality assessment approaches:

\begin{table}[h]
\centering
\caption{Source Frameworks for TERAN}
\label{tab:sources}
\begin{tabular}{p{3.5cm}p{4cm}p{3cm}p{3cm}}
\toprule
\textbf{Source} & \textbf{Framework} & \textbf{Contribution} & \textbf{Status} \\
\midrule
Belcher et al.\ (2016) & QAF & Primary structure & Peer-reviewed, 500+ cites \\
TACT (2023) & Trustworthiness & Validation & CFA-validated \\
Glassick et al.\ (1997) & Scholarship Assessed & Rigor criteria & Classic standard \\
Guba (1981) & Naturalistic Inquiry & Credibility & 2000+ citations \\
IDRC (2017) & Research Quality Plus & Development focus & Institutional \\
\bottomrule
\end{tabular}
\end{table}

\tagEMP{Belcher et al.\ 2016: Systematic review of 38 peer-reviewed articles on transdisciplinary research quality}

\subsection{The Belcher QAF Framework}

The primary source is Belcher et al.\ (2016), ``Defining and assessing research quality in a transdisciplinary context,'' published in \emph{Research Evaluation} (Oxford Academic). Their systematic review identified four overarching principles:

\begin{result}[Belcher QAF Four Principles]
\label{res:belcher}
\begin{enumerate}
    \item \textbf{Relevance:} Importance, significance, and usefulness to problem context
    \item \textbf{Credibility:} Robust findings, dependable knowledge sources
    \item \textbf{Legitimacy:} Fair, ethical, inclusive process
    \item \textbf{Effectiveness:} Generates knowledge, stimulates action
\end{enumerate}
\end{result}

\subsection{Mapping: Belcher QAF $\rightarrow$ BCM TERAN}

We adapt the four Belcher principles into five BCM-appropriate dimensions:

\begin{table}[h]
\centering
\caption{Mapping Belcher QAF to TERAN}
\label{tab:mapping}
\begin{tabular}{lll}
\toprule
\textbf{Belcher QAF} & \textbf{BCM TERAN} & \textbf{Rationale} \\
\midrule
Credibility (preparation) & \textbf{T}heory (25\%) & Literature grounding \\
Credibility (data/sources) & \textbf{E}vidence (30\%) & Empirical basis \\
Credibility (methods) & \textbf{R}igor (15\%) & Formal correctness \\
Relevance + Effectiveness & \textbf{A}pplicability (15\%) & Practical utility \\
Legitimacy & \textbf{N} Transparency (15\%) & Epistemic honesty \\
\bottomrule
\end{tabular}
\end{table}

\tagTHR{Mapping derived from Belcher Table 3 criterion definitions}

\subsection{Why Five Dimensions Instead of Four?}

Belcher's ``Credibility'' encompasses three distinct aspects that warrant separate assessment in a quantitative-parametric framework like EBF:

\begin{enumerate}
    \item \textbf{Theoretical preparation} (T): Are the right theories cited?
    \item \textbf{Empirical grounding} (E): Are parameters validated?
    \item \textbf{Methodological rigor} (R): Are derivations correct?
\end{enumerate}

\tagTHR{Decomposition justified by EBF's quantitative parameter claims}

\subsection{Weight Justification}

\begin{axiom}[Evidence Overweighting]
\label{ax:evidence}
For frameworks making quantitative parameter claims, Evidence (E) receives elevated weight (30\% vs.\ standard 25\%) because:
\begin{enumerate}
    \item Parameter values require empirical validation
    \item Unvalidated parameters propagate through all calculations
    \item Belcher emphasizes ``dependable knowledge sources''
\end{enumerate}
\end{axiom}

% ═══════════════════════════════════════════════════════════════════════════════
% SECTION 2: THE FIVE TERAN DIMENSIONS
% ═══════════════════════════════════════════════════════════════════════════════

\section{The Five TERAN Dimensions}
\label{sec:dimensions}

\subsection{T: Theory (25\%)}

\textbf{Definition:} The degree to which claims are grounded in established theory and literature.

\begin{table}[h]
\centering
\caption{Theory Subdimensions}
\begin{tabular}{clcp{6cm}}
\toprule
\textbf{Sub} & \textbf{Description} & \textbf{Max} & \textbf{Belcher Equivalent} \\
\midrule
T1 & Literature coverage & 2 & ``Broad preparation'' \\
T2 & Axiomatic consistency & 2 & ``Research approach fits purpose'' \\
T3 & Completeness & 1 & ``Broad preparation'' \\
T4 & Minimality & 1 & \textit{(BCM-specific)} \\
T5 & Proofs provided & 2 & ``Clearly presented argument'' \\
T6 & Novelty clarity & 2 & ``Contribution to knowledge'' \\
\midrule
& \textbf{Total} & \textbf{10} & \\
\bottomrule
\end{tabular}
\end{table}

\subsection{E: Evidence (30\%)}

\textbf{Definition:} The empirical basis for claims and parameters.

\begin{table}[h]
\centering
\caption{Evidence Subdimensions}
\begin{tabular}{clcp{6cm}}
\toprule
\textbf{Sub} & \textbf{Description} & \textbf{Max} & \textbf{Belcher Equivalent} \\
\midrule
E1 & Data sources used & 3 & ``Appropriate methods'' \\
E2 & Parameter calibration & 3 & ``Research approach fits purpose'' \\
E3 & Validation performed & 2 & ``Transferability'' \\
E4 & Falsifiability & 1 & ``Limitations stated'' \\
E5 & Robustness & 1 & ``Ongoing monitoring'' \\
\midrule
& \textbf{Total} & \textbf{10} & \\
\bottomrule
\end{tabular}
\end{table}

\textbf{Special Rule:} If a new construct (e.g., FEPSDE) lacks confirmatory factor analysis, maximum E score is 4/10.

\subsection{R: Rigor (15\%)}

\textbf{Definition:} Formal correctness of notation, definitions, and derivations.

\begin{table}[h]
\centering
\caption{Rigor Subdimensions}
\begin{tabular}{clcp{6cm}}
\toprule
\textbf{Sub} & \textbf{Description} & \textbf{Max} & \textbf{Glassick Equivalent} \\
\midrule
R1 & Notation consistency & 2 & ``Effective Presentation'' \\
R2 & Definition clarity & 2 & ``Clear Goals'' \\
R3 & Proof completeness & 3 & ``Appropriate Methods'' \\
R4 & Mathematical correctness & 3 & ``Appropriate Methods'' \\
\midrule
& \textbf{Total} & \textbf{10} & \\
\bottomrule
\end{tabular}
\end{table}

\subsection{A: Applicability (15\%)}

\textbf{Definition:} Practical usefulness and implementability.

\begin{table}[h]
\centering
\caption{Applicability Subdimensions}
\begin{tabular}{clcp{6cm}}
\toprule
\textbf{Sub} & \textbf{Description} & \textbf{Max} & \textbf{Belcher Equivalent} \\
\midrule
A1 & Worked examples & 3 & ``Practical application'' \\
A2 & Implementability & 3 & ``Practical application'' \\
A3 & Scalability & 2 & ``Transferability'' \\
A4 & Data requirements & 2 & ``Feasible research project'' \\
\midrule
& \textbf{Total} & \textbf{10} & \\
\bottomrule
\end{tabular}
\end{table}

\subsection{N: Transparency (15\%)}

\textbf{Definition:} Epistemic honesty about the status of claims.

\begin{table}[h]
\centering
\caption{Transparency Subdimensions}
\begin{tabular}{clcp{6cm}}
\toprule
\textbf{Sub} & \textbf{Description} & \textbf{Max} & \textbf{Guba Equivalent} \\
\midrule
N1 & Epistemic tags used & 3 & ``Confirmability'' \\
N2 & Established/novel separated & 3 & ``Confirmability'' \\
N3 & Peer-review status noted & 2 & \textit{(BCM-specific)} \\
N4 & Limitations discussed & 2 & ``Neutrality'' \\
\midrule
& \textbf{Total} & \textbf{10} & \\
\bottomrule
\end{tabular}
\end{table}

% ═══════════════════════════════════════════════════════════════════════════════
% SECTION 3: EPISTEMIC STATUS TAGS
% ═══════════════════════════════════════════════════════════════════════════════

\section{Epistemic Status Tags}
\label{sec:tags}

\subsection{The Five-Level System}

Every parameter, claim, and construct in EBF should carry an epistemic tag indicating its evidential basis:

\begin{table}[h]
\centering
\caption{Epistemic Status Tags}
\label{tab:tags}
\begin{tabular}{ccp{7cm}c}
\toprule
\textbf{Tag} & \textbf{Level} & \textbf{Definition} & \textbf{Color} \\
\midrule
\tagEMP{} & $\star\star\star\star\star$ & Empirically validated in peer-reviewed publication & Green \\
\tagTHR{} & $\star\star\star\star\phantom{\star}$ & Theoretically derived from established axioms & Blue \\
\tagLLM{} & $\star\star\star\phantom{\star\star}$ & Estimated via LLM-Monte-Carlo simulation & Yellow \\
\tagILL{} & $\star\star\phantom{\star\star\star}$ & Illustrative example value only & Orange \\
\tagHYP{} & $\star\phantom{\star\star\star\star}$ & Hypothetical/speculative & Red \\
\bottomrule
\end{tabular}
\end{table}

\subsection{Tag Syntax}

\textbf{LaTeX:}
\begin{verbatim}
$\gamma_{FP} = 0.09$ \tagEMP{Tversky \& Kahneman 1992}
$\delta_S = 0.04$ \tagLLM{n=500, Claude 3.5, Jan 2026}
$\beta = 0.6$ \tagILL{Anna persona illustration}
\end{verbatim}

\subsection{Aggregation Rule}

\begin{axiom}[Weakest Link Aggregation]
When a construct depends on multiple parameters, its aggregate epistemic status equals the \emph{weakest} component tag:
\begin{center}
\begin{tabular}{ll}
All EMP & $\rightarrow$ EMP \\
EMP + THR & $\rightarrow$ THR \\
Contains LLM & $\rightarrow$ LLM \\
Contains ILL & $\rightarrow$ ILL \\
Contains HYP & $\rightarrow$ HYP \\
\end{tabular}
\end{center}
\end{axiom}

\subsection{Tag Coverage Metric}

\begin{definition}[Tag Coverage]
\begin{equation}
\text{Tag Coverage} = \frac{\text{Parameters with tags}}{\text{Total parameters}} \times 100\%
\end{equation}
\textbf{Target:} $\geq 90\%$ for publication readiness.
\end{definition}

% ═══════════════════════════════════════════════════════════════════════════════
% SECTION 4: APPLICATION
% ═══════════════════════════════════════════════════════════════════════════════

\section{Application: Worked Example}
\label{sec:application}

\subsection{Initial Assessment: 5C CORE (2026-01-06)}

\begin{table}[h]
\centering
\caption{TERAN Scores for 5C CORE Appendices}
\label{tab:assessment}
\begin{tabular}{lccccccc}
\toprule
\textbf{Appendix} & \textbf{T} & \textbf{E} & \textbf{R} & \textbf{A} & \textbf{N} & \textbf{Total} & \textbf{Stars} \\
\midrule
AAA (WHO) & 6 & 2 & 6 & 4 & 5 & 4.3 & $\star\star$ \\
B (HOW) & 9 & 5 & 7 & 6 & 5 & 6.4 & $\star\star\star$ \\
C (WHAT) & 6 & 3 & 6 & 6 & 4 & 4.7 & $\star\star$ \\
V (WHEN) & 7 & 5 & 6 & 6 & 6 & 5.9 & $\star\star\star$ \\
BBB (WHERE) & 7 & 5 & 7 & 6 & 7 & 6.2 & $\star\star\star$ \\
\midrule
\textbf{Average} & 7.0 & 4.0 & 6.4 & 5.6 & 5.4 & \textbf{5.5} & $\star\star\star$ \\
\bottomrule
\end{tabular}
\end{table}

\textbf{Key Finding:} Evidence (E = 4.0) is the weakest dimension, driven by:
\begin{itemize}
    \item FEPSDE lacks confirmatory factor analysis
    \item $\gamma$-coefficients are illustrative, not calibrated
    \item Level structure (AAA) has no empirical validation
\end{itemize}

\subsection{Tag Coverage Audit}

\begin{table}[h]
\centering
\caption{Epistemic Tag Coverage by Appendix}
\begin{tabular}{lccc}
\toprule
\textbf{Appendix} & \textbf{Parameters} & \textbf{Tagged} & \textbf{Coverage} \\
\midrule
AAA & $\sim$20 & 5 & 25\% \\
B & $\sim$163 & 10 & 6\% \\
C & $\sim$144 & 5 & 3\% \\
V & $\sim$50 & 15 & 30\% \\
BBB & $\sim$100 & 40 & 40\% \\
\midrule
\textbf{Total} & $\sim$477 & $\sim$75 & \textbf{16\%} \\
\bottomrule
\end{tabular}
\end{table}

\textbf{Critical Gap:} 84\% of parameters lack epistemic status tags.

% ═══════════════════════════════════════════════════════════════════════════════
% SECTION 5: CRITICAL FOUNDATIONS
% ═══════════════════════════════════════════════════════════════════════════════

\section{Critical Foundations}

\subsection{Why Not Use Belcher QAF Directly?}

\textbf{Objection:} ``Why adapt rather than adopt Belcher's framework as-is?''

\textbf{Response:}
\begin{enumerate}
    \item \textbf{Domain specificity:} EBF is a quantitative-parametric framework; Belcher QAF was designed for transdisciplinary research broadly
    \item \textbf{Parameter focus:} EBF makes specific numerical claims ($\gamma_{FP} = 0.09$) requiring explicit calibration assessment
    \item \textbf{Epistemic transparency:} The tag system (EMP/THR/LLM/ILL/HYP) is not part of Belcher but essential for EBF
    \item \textbf{Credibility decomposition:} Separating T/E/R provides finer diagnostic resolution
\end{enumerate}

\subsection{Is TERAN Validated?}

\textbf{Objection:} ``TERAN itself lacks empirical validation.''

\textbf{Response:}
\begin{itemize}
    \item The \emph{dimensions} derive from validated frameworks (Belcher: systematic review; TACT: CFA-validated)
    \item The \emph{weights} are theoretically justified, not empirically optimized
    \item TERAN is a \emph{diagnostic tool}, not a predictive model—validation is through utility, not prediction accuracy
    \item Status: \tagTHR{Derived from peer-reviewed frameworks}
\end{itemize}

\subsection{Why Weight Evidence at 30\%?}

\textbf{Objection:} ``Evidence weighting seems arbitrary.''

\textbf{Response:}
\begin{itemize}
    \item Standard equal weighting (20\% each) would underweight the critical gap in EBF
    \item EBF makes quantitative claims that propagate through all calculations
    \item Belcher et al.\ emphasize ``dependable knowledge sources'' as central to credibility
    \item Sensitivity analysis: Even at 25\%, Evidence remains the weakest dimension
\end{itemize}

% ═══════════════════════════════════════════════════════════════════════════════
% SECTION 6: OPEN ISSUES
% ═══════════════════════════════════════════════════════════════════════════════


%%%%%%%%%%%%%%%%%%%%%%%%%%%%%%%%%%%%%%%%%%%%%%%%%%%%%%%%%%%%%%%%%%%%%%%%%%%%%%%
\section{References}
\label{sec:references}
%%%%%%%%%%%%%%%%%%%%%%%%%%%%%%%%%%%%%%%%%%%%%%%%%%%%%%%%%%%%%%%%%%%%%%%%%%%%%%%

See master bibliography (\texttt{bcm\_master.bib}) for complete citations.

\nocite{bcm_master}


%%%%%%%%%%%%%%%%%%%%%%%%%%%%%%%%%%%%%%%%%%%%%%%%%%%%%%%%%%%%%%%%%%%%%%%%%%%%%%%
\section{Glossary of Symbols}
\label{sec:qly-glossary}
%%%%%%%%%%%%%%%%%%%%%%%%%%%%%%%%%%%%%%%%%%%%%%%%%%%%%%%%%%%%%%%%%%%%%%%%%%%%%%%

For comprehensive definitions, see Appendix G (Master Glossary).

\begin{table}[htbp]
\centering
\caption{Key Symbols in Appendix UNMAPPED_QLY}
\begin{tabular}{lll}
\toprule
\textbf{Symbol} & \textbf{Meaning} & \textbf{Reference} \\
\midrule
$\gamma$ & Complementarity parameter & Appendix UNMAPPED_B \\
$\Psi$ & Context vector & Appendix UNMAPPED_V \\
$U$ & Utility function & Chapter 3 \\
\bottomrule
\end{tabular}
\end{table}


\section{Open Issues}

\subsection{Methodological}

\begin{enumerate}
    \item \textbf{Inter-rater reliability:} TERAN assessments may vary across assessors; calibration studies needed
    \item \textbf{Weight sensitivity:} How robust are conclusions to alternative weightings?
    \item \textbf{Automated assessment:} Can tag coverage be computed automatically from LaTeX source?
\end{enumerate}

\subsection{Theoretical}

\begin{enumerate}
    \item \textbf{LLM tag validity:} What confidence level justifies LLM-MC estimates?
    \item \textbf{Tag aggregation:} Is weakest-link appropriate, or should weighted averaging be used?
    \item \textbf{Publication threshold:} Is 8.0 the right threshold for ``publication ready''?
\end{enumerate}

\subsection{Empirical}

\begin{enumerate}
    \item \textbf{Benchmark validation:} Compare TERAN scores to actual publication outcomes
    \item \textbf{Expert calibration:} Do domain experts agree with TERAN assessments?
    \item \textbf{Longitudinal tracking:} Does quality improve as predicted by TERAN diagnostics?
\end{enumerate}

% ═══════════════════════════════════════════════════════════════════════════════
% SECTION 7: SUMMARY
% ═══════════════════════════════════════════════════════════════════════════════


%%%%%%%%%%%%%%%%%%%%%%%%%%%%%%%%%%%%%%%%%%%%%%%%%%%%%%%%%%%%%%%%%%%%%%%%%%%%%%%
\section{Key Results}
\label{sec:qly-results}
%%%%%%%%%%%%%%%%%%%%%%%%%%%%%%%%%%%%%%%%%%%%%%%%%%%%%%%%%%%%%%%%%%%%%%%%%%%%%%%

The analysis yields the following key results:

\begin{enumerate}
    \item \textbf{Result 1:} [Primary finding with quantitative support]
    \item \textbf{Result 2:} [Secondary finding with implications]
    \item \textbf{Result 3:} [Practical application or boundary condition]
\end{enumerate}


\section{Summary}

\begin{tcolorbox}[colback=gray!10,colframe=gray!50!black,title=\textbf{Key Results}]
\begin{enumerate}
    \item \textbf{TERAN Framework:} Five dimensions (Theory, Evidence, Rigor, Applicability, Transparency) adapted from Belcher et al.\ (2016) QAF
    
    \item \textbf{Epistemic Tags:} Five-level system (EMP/THR/LLM/ILL/HYP) for transparent parameter provenance
    
    \item \textbf{Current Status:} 5C CORE averages 5.5/10 (Working Paper level); Evidence is weakest at 4.0/10
    
    \item \textbf{Critical Gap:} 84\% of parameters lack epistemic status tags
    
    \item \textbf{Target Trajectory:} 5.5 (Jan 2026) $\rightarrow$ 7.0 (Q2 2026) $\rightarrow$ 8.0 (Q4 2026)
\end{enumerate}
\end{tcolorbox}

% ═══════════════════════════════════════════════════════════════════════════════
% GLOSSARY LINK
% ═══════════════════════════════════════════════════════════════════════════════

\section*{Glossary}

Key terms defined in Appendix UNMAPPED_GLS (Master Glossary):
\begin{itemize}
    \item \textbf{TERAN:} Theory-Evidence-Rigor-Applicability-Transparency assessment framework
    \item \textbf{Epistemic Status:} The evidential basis for a claim or parameter
    \item \textbf{Tag Coverage:} Percentage of parameters with explicit epistemic tags
    \item \textbf{Template Compliance:} Adherence to formal structural requirements
    \item \textbf{Content Quality:} Scientific soundness and empirical grounding
\end{itemize}

$\rightarrow$ \textbf{Full Glossary:} See Appendix UNMAPPED_GLS

% ═══════════════════════════════════════════════════════════════════════════════
% REFERENCES
% ═══════════════════════════════════════════════════════════════════════════════

\section*{References}

\begin{enumerate}
    \item \textbf{Belcher, B.M., Rasmussen, K.E., Kemshaw, M.R., \& Zornes, D.A.} (2016). Defining and assessing research quality in a transdisciplinary context. \emph{Research Evaluation}, 25(1), 1--17. \tagEMP{Primary source, 500+ citations}
    
    \item \textbf{Daniel, B.K.} (2018). Empirically grounded framework for rigor in qualitative research. \emph{Foundation for TACT}. \tagEMP{CFA-validated 2023}
    
    \item \textbf{Glassick, C.E., Huber, M.T., \& Maeroff, G.I.} (1997). \emph{Scholarship Assessed: Evaluation of the Professoriate}. Jossey-Bass. \tagEMP{Classic standard}
    
    \item \textbf{Guba, E.G.} (1981). Criteria for assessing the trustworthiness of naturalistic inquiries. \emph{ECTJ}, 29, 75--91. \tagEMP{2000+ citations}
    
    \item \textbf{IDRC} (2017). Research Quality Plus (RQ+) Assessment Instrument. International Development Research Centre. \tagEMP{Institutional standard}
    
    \item \textbf{Cash, D.W., et al.} (2003). Knowledge systems for sustainable development. \emph{PNAS}, 100(14), 8086--8091. \tagEMP{Credibility-Relevance-Legitimacy framework}
\end{enumerate}

$\rightarrow$ \textbf{Full Bibliography:} See Appendix THL1 (Master Bibliography)

% ═══════════════════════════════════════════════════════════════════════════════
% END DOCUMENT
% ═══════════════════════════════════════════════════════════════════════════════


% Link to master bibliography
\nocite{bcm_master}

\end{document}
